% COLLECTION_METADATA: {"schema_version": 1, "kind": "reused_writeup", "independently_reviewed": false, "primary_source": "attacks/erdos/gpt_pro_5.2/107.tex", "source_paths": ["attacks/erdos/gpt_5.2/107.tex", "attacks/erdos/gpt_pro_5.2/107.tex"], "source_model": "gpt_pro_5.2", "source_completion": null, "source_claim": "unresolved"}
\section*{Erd\H{o}s Problem 107 --- collection record}

\subsection*{PROVENANCE AND STATUS}
Official problem: https://www.erdosproblems.com/107

Saved collaborative-database status: falsifiable.
Snapshot checked: 2026-09-23T17:22:24Z.
Latest status and formalization links: https://teorth.github.io/erdosproblems/

This is an attributed collection of existing writeups. The model-folder names below identify their original attribution; placement in GPT\_6\_Astra\_Ultra does not attribute their mathematical work to that model. All locally available source versions selected for this problem are reproduced below, without editing their text.

Proof claims, computations, literature statements, and completion estimates in the reused text have not been independently verified in this collection pass. A source claiming a solution does not change the saved upstream status.

Primary source (latest numbered version in the preferred model folder): \texttt{attacks/erdos/gpt\_pro\_5.2/107.tex}.

\subsection*{REUSED SOURCE 1}
Original attribution: \texttt{gpt\_5.2}.
Original file: \texttt{attacks/erdos/gpt\_5.2/107.tex}.

% BEGIN REUSED SOURCE: attacks/erdos/gpt_5.2/107.tex

OUTPUT
1) FORMAL RESTATEMENT
Definitions/Conventions.
- A finite set $P\subset\mathbb R^2$ is in \emph{general position} if no three points are collinear.
- For $n\ge 3$, a subset $Q\subseteq P$ of size $n$ \emph{forms the vertices of a convex $n$-gon} if the $n$ points of $Q$ are in convex position, i.e. every point of $Q$ is a vertex of $\mathrm{conv}(Q)$.
- Define $f(n)$ to be the least integer $N$ such that every set $P\subset\mathbb R^2$ of $N$ points in general position contains an $n$-point subset in convex position.

Conjecture (Erd\H{o}s--Szekeres).
Prove that for all $n\ge 3$,
\[f(n)=2^{n-2}+1.\]

2) QUICK LITERATURE/CONTEXT CHECK
Web browsing is not available in this session.
From the problem statement:
- The exact values $f(4)=5$ and $f(5)=9$ are known.
- Erd\H{o}s--Szekeres proved $2^{n-2}+1\le f(n)\le \binom{2n-4}{n-2}+1$.
- Stronger upper bounds are known, but the exact formula is open in general.
I do not use the deep upper bounds here.

3) ATTACK PLAN
Proof track ideas:
1. Try to prove the conjectured exact formula by understanding extremal configurations with no convex $n$-gon and showing they have size at most $2^{n-2}$.
2. For small $n$, prove exact values directly as a reality check and as a base for any induction attempt.

Disproof track ideas:
1. Try to construct, for some $n$, a set of size $2^{n-2}+1$ in general position with no convex $n$-gon. Any such construction would refute the conjecture.

Chosen path in this attempt: prove two basic, fully rigorous facts: monotonicity of $f(n)$, and the exact value $f(4)=5$ (the first nontrivial case).

4) WORK
Lemma 1 (Monotonicity).
For all $n\ge 3$, one has $f(n+1)\ge f(n)$.

Proof.
Let $N:=f(n+1)$. By definition, every set of $N$ points in general position contains a subset of $n+1$ points in convex position.
Any subset of vertices of a convex polygon is again in convex position (because removing vertices cannot create an interior point among the remaining ones).
Therefore, the same $N$ points also contain $n$ points in convex position.
Thus $N$ is admissible for $f(n)$, so $f(n)\le N=f(n+1)$. \qed

Lemma 2 (Exact value $f(4)=5$).
Every set of $5$ points in $\mathbb R^2$ in general position contains $4$ points in convex position, and there exists a set of $4$ points in general position with no $4$ points in convex position.
Hence $f(4)=5$.

Proof.
Lower bound $f(4)\ge 5$: take three noncollinear points $A,B,C$ and a fourth point $D$ strictly inside the triangle $\triangle ABC$.
Then $A,B,C,D$ are in general position, but $D$ is not a vertex of the convex hull of $\{A,B,C,D\}$, so these four points are not in convex position.
Thus $f(4)>4$, i.e. $f(4)\ge 5$.

Upper bound $f(4)\le 5$: let $P$ be any set of $5$ points in general position.
If at least $4$ points of $P$ lie on the boundary of $\mathrm{conv}(P)$, then those $4$ boundary points are in convex position and we are done.
So assume $\mathrm{conv}(P)$ has at most $3$ vertices, meaning at least one point of $P$ lies strictly inside the convex hull of the others.

Consider any straight-line triangulation of the point set $P$ (a planar subdivision of $\mathrm{conv}(P)$ into triangles whose vertices are points of $P$ and whose edges are noncrossing segments between points of $P$).
Such a triangulation exists for any finite set of points in general position.
Because $P$ is not all on the boundary of its convex hull, the triangulation has at least one \emph{interior} edge $ab$ (an edge not on the boundary of $\mathrm{conv}(P)$).
This interior edge is incident to exactly two triangles in the triangulation, say triangles $abc$ and $abd$ with $c\ne d$.

Since $ab$ is an interior edge, the triangles $abc$ and $abd$ lie on opposite sides of the line through $a$ and $b$.
Equivalently, the points $c$ and $d$ lie in opposite open half-planes determined by the line $ab$.
In particular, the segment $\overline{cd}$ intersects the segment $\overline{ab}$ at an interior point of both segments.
Whenever two segments $\overline{ab}$ and $\overline{cd}$ cross (and no three of $a,b,c,d$ are collinear), the four endpoints $a,b,c,d$ are in convex position: they are the vertices of a convex quadrilateral with diagonals $\overline{ab}$ and $\overline{cd}$.
Thus $\{a,b,c,d\}\subseteq P$ is a 4-point subset in convex position.

This proves every $5$-point general-position set contains a convex quadrilateral, i.e. $f(4)\le 5$.

Combining both bounds yields $f(4)=5$. \qed

FAST REALITY CHECK
- The conjectured formula gives $f(3)=2^{1}+1=3$ and $f(4)=2^{2}+1=5$, matching the trivial triangle case and Lemma 2.
- The statement reports $f(5)=2^{3}+1=9$, consistent with the conjecture but not proved here.

5) VERIFICATION
- Lemma 1 uses only the fact that subsets of vertices of a convex polygon remain in convex position.
- In Lemma 2, the key geometric step is: if a line through two interior points intersects two sides of a triangle, then two triangle vertices lie on opposite sides of that line; then crossing diagonals imply convex position of four points. No collinearity assumptions ensure the crossing criterion is valid.

6) FINAL
**UNRESOLVED**

(i) Strongest fully proved partial result:
I proved monotonicity of $f(n)$ (Lemma 1) and established the exact value $f(4)=5$ with a complete proof (Lemma 2).

(ii) First gap:
Prove $f(n)\le 2^{n-2}+1$ for general $n$ (or find a counterexample for some $n$).
Beyond $n=5$, known proofs require substantially deeper combinatorial geometry than the small-case argument in Lemma 2.

(iii) Top 3 next moves (concrete):
1. Prove the next open-looking case(s) directly (e.g. re-derive $f(5)=9$) as a stepping stone and to understand extremal configurations.
2. Study the structure of known lower-bound constructions of size $2^{n-2}$ with no convex $n$-gon (as mentioned in the statement) and attempt to show they are extremal.
3. Try to connect the conjecture to an induction on $n$ by identifying a canonical ``splitting vertex'' in a point set that reduces the problem to smaller $n$.

(iv) What a minimal counterexample would likely look like (structure/parameters):
For some $n\ge 6$, a set of $2^{n-2}+1$ points in general position with no $n$ points in convex position; such a configuration would have to be highly structured and likely resemble known extremal constructions for smaller $n$.

% END REUSED SOURCE: attacks/erdos/gpt_5.2/107.tex

\subsection*{REUSED SOURCE 2}
Original attribution: \texttt{gpt\_pro\_5.2}.
Original file: \texttt{attacks/erdos/gpt\_pro\_5.2/107.tex}.

% BEGIN REUSED SOURCE: attacks/erdos/gpt_pro_5.2/107.tex

\section*{Erd\H{o}s Problem \#107}

\subsection*{1) FORMAL RESTATEMENT}
For an integer $n\ge3$, let $f(n)$ be the least integer $N$ such that every set $P\subset\mathbb R^2$ of $N$ points with \emph{no three collinear} contains $n$ points in \emph{convex position}, i.e. $n$ points that are the vertices of a convex $n$-gon (equivalently, no one of the $n$ points lies in the convex hull of the other $n-1$).

The conjectured exact value is
\[
 f(n)=2^{n-2}+1\quad(n\ge3).
\]

\subsection*{2) QUICK LITERATURE/CONTEXT CHECK}
I only record what the problem text states:
\begin{itemize}
\item Erd\H{o}s--Szekeres proved the bounds $2^{n-2}+1\le f(n)\le \binom{2n-4}{n-2}+1$.
\item Several improvements exist; the best stated bound is $f(n)\le 2^{n+O(\sqrt{n\log n})}$.
\end{itemize}
No other results are assumed/proved here.

\subsection*{3) ATTACK PLAN}
\begin{itemize}
\item \textbf{Proof track (full conjecture):} show every set of $2^{n-2}+1$ points contains a convex $n$-gon, likely via a sharp recursion or a refined partial-order argument.
\item \textbf{Disproof track:} exhibit, for some $n$, a configuration of $2^{n-2}+1$ points with no convex $n$-gon.
\end{itemize}
In this writeup I give a complete proof of the base case $f(4)=5$ and basic structural lemmas, but I do not resolve the general conjecture.

\subsection*{4) WORK}
\paragraph{Lemma 107.1 (the case $n=4$: $f(4)=5$).}
One has $f(4)=5$.

\emph{Proof.}
\emph{Lower bound $f(4)\ge5$.}  There exists a $4$-point set with no convex quadrilateral: take three vertices of a triangle and a fourth point strictly inside that triangle.  Any four-point subset is the whole set, and one point lies in the convex hull of the other three, so there is no convex $4$-gon.  Thus $f(4)\ge5$.

\emph{Upper bound $f(4)\le5$.}  Let $P$ be any set of $5$ points with no three collinear.  Consider all triangles formed by triples of points in $P$, and choose a triangle $\triangle ABC$ of \emph{minimum (Euclidean) area}.

\emph{Claim:} no other point of $P$ lies strictly inside $\triangle ABC$.
Indeed, if some point $D\in P$ lies inside $\triangle ABC$, then the three triangles $\triangle ABD$, $\triangle BCD$, $\triangle CAD$ form a partition of $\triangle ABC$ into three nondegenerate triangles.  In particular, each of these triangles has area strictly smaller than $\triangle ABC$, contradicting the minimality of $\triangle ABC$.

Therefore, the remaining two points of $P\setminus\{A,B,C\}$ lie \emph{outside} the triangle $\triangle ABC$.  Hence the convex hull $\operatorname{conv}(P)$ cannot be a triangle (a triangular convex hull would contain all points in or on that triangle, with at least two points inside, contradicting the claim).  Thus $\operatorname{conv}(P)$ has at least $4$ vertices, and those $4$ vertices form a convex quadrilateral.

So every $5$-point set contains a convex $4$-gon, i.e. $f(4)\le5$.  Together with the lower bound, $f(4)=5$. \qed

\paragraph{Lemma 107.2 (monotonicity and trivial bounds).}
For $n\ge3$:
\begin{enumerate}
\item $f(n+1)\ge f(n)$.
\item $f(n)\ge n$.
\end{enumerate}

\emph{Proof.}
(2) is immediate since fewer than $n$ points cannot contain $n$ points in convex position.

For (1): by definition, any set of $f(n+1)$ points in general position contains $n+1$ points in convex position.  Removing any one of these $n+1$ vertices leaves $n$ points still in convex position.  Hence any set of $f(n+1)$ points contains a convex $n$-gon, so $f(n)\le f(n+1)$, i.e. $f(n+1)\ge f(n)$. \qed

\subsection*{5) VERIFICATION (FAST REALITY CHECK)}
\begin{itemize}
\item $f(3)=3$: any $3$ noncollinear points form a triangle.
\item Lemma~107.1 gives $f(4)=5$.
\item The conjecture predicts $f(5)=2^{3}+1=9$, consistent with the statement in the problem text that Tur\'an and Makai proved $f(5)=9$.
\end{itemize}

\subsection*{6) FINAL}
\textbf{UNRESOLVED.}

(i) \emph{Strongest fully proved partial result obtained here.}
A complete proof that $f(4)=5$ (Lemma~107.1), plus basic monotonicity/trivial bounds (Lemma~107.2).

(ii) \emph{Exact first gap.}
Proving (or disproving) the general identity $f(n)=2^{n-2}+1$ for all $n\ge3$.

(iii) \emph{Top 3 next moves (concrete targets).}
\begin{enumerate}
\item Re-prove the general Erd\H{o}s--Szekeres upper bound $f(n)\le \binom{2n-4}{n-2}+1$ in a fully self-contained way as a stepping stone.
\item Attempt to sharpen the induction/recursion to match the lower bound $2^{n-2}+1$.
\item Computationally search for extremal configurations for $n=6,7$ to test structure and potential stability.
\end{enumerate}

(iv) \emph{Minimal counterexample structure.}
A disproof would require some $n$ and an explicit set of $2^{n-2}+1$ points in general position containing no convex $n$-gon; it would have to beat the classical Erd\H{o}s--Szekeres construction of size $2^{n-2}$ without creating a convex $n$-tuple.



% END REUSED SOURCE: attacks/erdos/gpt_pro_5.2/107.tex

\subsection*{COMPLETION ESTIMATE}
The primary source supplies no extractable numerical completion estimate. Its mathematical progress is not numerically assessed here.

Newly verified progress in this collection pass: $0\%$. This measures new verification work, not the amount of mathematics achieved in the reused text.
